\begin{abstract}
    As an emerging technology that has attracted huge attention, non-line-of-sight (NLOS) imaging can reconstruct hidden objects by analyzing the diffuse reflection on a relay surface, with broad application prospects in the fields of autonomous driving, medical imaging, and defense. Despite the challenges of low signal-to-noise ratio (SNR) and high ill-posedness, NLOS imaging has been developed rapidly in recent years. Most current NLOS imaging technologies use conventional physical models, constructing imaging models through active or passive illumination and using reconstruction algorithms to restore hidden scenes. Moreover, deep learning algorithms for NLOS imaging have also received much attention recently. This paper presents a comprehensive overview of both conventional and deep learning-based NLOS imaging techniques. It is an updated version of our peer-reviewed journal paper published in \emph{APSIPA Transactions on Signal and Information Processing} (2022, vol.~11, pp.~1--48, \href{https://doi.org/10.1561/116.00000019}{DOI:~10.1561/116.00000019}), extended to cover significant advances from 2022 to 2026; \textbf{the content added in this update has not undergone peer review}. In particular, we highlight the emergence of transformer-based and diffusion-model-based reconstruction networks, the achievement of real-time dynamic NLOS video, the extension of NLOS imaging to non-optical modalities such as radar and ultrasound, and the exploration of new application scenarios including human pose estimation and super-field-of-view imaging. Besides, we survey new proposed NLOS scenes, and discuss the challenges and prospects of existing technologies. Such a survey can help readers have an overview of different types of NLOS imaging, thus expediting the development of seeing around corners. A curated list of 150+ NLOS papers and an interactive project homepage are available at \url{https://github.com/ruixv/NLOS_Overview} and \url{https://ruixv.github.io/NLOS_Overview/}.
\end{abstract}
